\SourceLeaf{312}%
and the numbers given carry the expansion through quantities of eighth order inclusive. Thus \(t,l,b\) themselves are determined through quantities of ninth order inclusive, or up to but excluding eleventh order.

The expansion of \(\log b\) can also be carried out in another way, in powers of \(\beta,\tau\) and \(r\). If we set
\[
z=\operatorname{tang}y
\]
we obtain
\[
y=z-\tfrac13z^3+\tfrac15z^5-\tfrac17z^7+\tfrac19z^9-\text{etc.}
\]
and from this
\[
\log y=\log z-\tfrac13zz+\tfrac{13}{90}z^4-\tfrac{251}{2835}z^6+\tfrac{3551}{56700}z^8-\text{etc.}
\]
Applying this series to the equation
\[
\operatorname{tang}\tfrac12b=\frac{\operatorname{tang}\frac12t}{\operatorname{tang}B\mathord{.}\operatorname{tang}T}
\]
after using the substitutions
\[
\begin{aligned}
\operatorname{tang}B\mathord{.}\operatorname{tang}T&=\frac\tau\beta\\
\operatorname{tang}\tfrac12t&=\frac{m\tau}{2\BareRoot{(1-\frac14mm\tau\tau)}}
\end{aligned}
\]
to put it in the following form
\[
\frac{\beta m}{2\BareRoot{(1-\frac14mm\tau\tau)}}=\operatorname{tang}\tfrac12b
\]
gives
\[
\begin{aligned}
\log b={}&\log\beta+\log m+\tfrac18mm\tau\tau+\tfrac1{64}m^4\tau^4+\tfrac1{384}m^6\tau^6+\tfrac1{2048}m^8\tau^8+\text{etc.}\\
&-\tfrac1{12}\beta\beta(mm+\tfrac14m^4\tau\tau+\tfrac1{16}m^6\tau^4+\tfrac1{64}m^8\tau^6+\text{etc.})\\
&+\tfrac{13}{1440}\beta^4(m^4+\tfrac12m^6\tau\tau+\tfrac3{16}m^8\tau^4+\text{etc.})\\
&-\tfrac{251}{181440}\beta^6(m^6+\tfrac34m^8\tau\tau+\text{etc.})\\
&+\tfrac{3551}{14515200}\beta^8(m^8+\text{etc.})\\
&+\text{etc.}
\end{aligned}
\]
or, collecting terms of the same order,
